\begin{textsample}{POS Dim 3 – rj – Score 17.00 – rj/rj\_inf\_transicoes\_2.txt}  \label{ex:f3_pos_020}
Da \textbf{Creche} para a Pré-escola Temos uma organização diferenciada em relação ao espaço e tempo na \textbf{creche} em comparação com a \textbf{pré-escola}.
Na \textbf{creche}, em sua maioria, as \textbf{crianças} passam pelo período integral, enquanto na \textbf{pré-escola} esse tempo é reduzido para parcial.
A \textbf{rotina} passa a ser modificada.
Requer muitas vezes que objetivos a serem alcançados pelas \textbf{crianças} na \textbf{creche} tenham sido atingidos, como, por exemplo, algumas ações que envolvam a autonomia, para uma melhor integração na \textbf{pré-escola}.
O que o \textbf{adulto} precisava fazer pela \textbf{criança}, a partir desta faixa \textbf{etária} de 4 anos ela já \textbf{consegue} realizar sozinha.
Esse período acaba causando certa insegurança nos pais e nas próprias \textbf{crianças}.
Em 2013, a Lei nº 12.796/2013 também estabeleceu que a Educação \textbf{Infantil} — contempla \textbf{crianças} de 4 e 5 anos na \textbf{pré-escola} — seria organizada com carga horária \textbf{mínima} anual de 800 horas, distribuída por, no mínimo, 200 \textbf{dias} letivos.
O atendimento à \textbf{criança} deve ser de, no mínimo, quatro horas \textbf{diárias} para o turno parcial e de sete horas para a jornada integral.
A norma já valia para o Ensino Fundamental e Médio.
Essa ampliação da obrigatoriedade do ensino encontra ressonância já na primeira meta do Plano Nacional de Educação – PNE ( BRASIL, 2014 ).
Ela trata da universalização da Educação \textbf{Infantil} na \textbf{pré-escola} – para as \textbf{crianças} de quatro e cinco anos de idade – e da ampliação da oferta de Educação \textbf{Infantil}, em \textbf{creches}, de forma a atender, no mínimo, 50 \% das \textbf{crianças} de até três anos, ao final da vigência do PNE, em 2024.
Os investimentos em Educação \textbf{Infantil} estão entre as prioridades, para garantir que todas as \textbf{crianças} tenham a oportunidade de receber os \textbf{cuidados} e os \textbf{estímulos} pedagógicos necessários.
Essa política educacional também compartilha a premissa de alguns estudos no sentido de que a \textbf{criança} que não passa pela Educação \textbf{Infantil} tem mais dificuldades na fase de alfabetização.
Isso se torna muita das vezes um problema, pois muitos educadores acreditam que alfabetizar as \textbf{crianças} na faixa \textbf{etária} da \textbf{pré-escola} resolveria os problemas da alfabetização e do Ensino Fundamental.
Deixam o \textbf{cuidar} - o educar e o \textbf{brincar} - em segundo plano, muitas vezes exigindo das \textbf{crianças} tarefas e atividades que comprometem o seu direito de ser \textbf{criança}.
Da Pré-escola para o Ensino Fundamental Muitas vezes essa trajetória da Educação \textbf{Infantil} ao Ensino Fundamental não é compreendida como um processo contínuo, mas precisamos deixar evidente que a Educação \textbf{Infantil} não é período preparatório para o Ensino Fundamental, e muito menos processo precoce de escolarização.
As Leis Federais 11.114/2005, 11.274/2006 e 12.796/2013 ( que altera a Lei 9.394/96 em acordo com a Emenda Constitucional 59/2009 ) instituíram uma nova organização do Ensino Fundamental.
Algumas delas definiram a antecipação do Ensino Fundamental, a ser iniciado aos 6 anos de idade, e ampliaram sua duração para nove anos.
Outras ampliaram a obrigatoriedade escolar, estendendo -a dos 4 aos 17 anos de idade.
Essa nova organização traz inevitáveis questionamentos, entre os quais os relacionados às etapas educativas.
Os documentos oficiais tratam dessa articulação, mas não explicam como ela pode ocorrer.
As DCNEI consistem no documento que tenta aproximar essas duas etapas.
Neste ponto temos várias pesquisas que apontam o impasse dessa transição.
De acordo com Batista e Neves ( 2016 ), precisamos da construção de práticas pedagógicas que integrem o educar e o \textbf{cuidar}, centradas nas interações e nas \textbf{brincadeiras}, que considerem a relação com as famílias, além da necessidade de diálogo com a continuidade do processo de escolarização da \textbf{infância}, considerando que o foco principal das práticas educativas, tanto na Educação \textbf{Infantil} quanto nos anos iniciais do Ensino Fundamental, é o sujeito e sua relação com a cultura.
Assim, seria possível a construção de uma continuidade educativa no processo de escolarização da \textbf{criança}.

% matched lemmas: adulto, brincadeira, brincar, conseguir, creche, criança, cuidado, cuidar, dia, diário, estímulo, etário, infantil, infância, pequeno, pré-escola, rotina
\end{textsample}
